\clearpage
\fancyhead[L]{\small DETERMINING LONGITUDE BY CHRONOMETER.}
\hypertarget{native-460-chronometric}{}
\expandafter\def\csname HistoricalNoteText462\endcsname{[observed by Dr. \textsc{Tiarks} during the expedition organized by the English Admiralty in 1824 to determine by chronometer the differences in longitude between Altona, Bremen, Helgoland and Greenwich. \textsc{Schering}.]}
\begin{center}{\large DETERMINING LONGITUDE BY CHRONOMETER.}\end{center}

\begin{center}
\rule{0.55\linewidth}{0.4pt}\\[2pt]
Astronomische Nachrichten. Vol. V. No. 110. cols. 227\,..\,234. November 1826.\\[2pt]
\rule{0.55\linewidth}{0.4pt}
\end{center}

Let $\theta, \theta', \theta''\ldots$ be the times ($n$ in all) when the chronometer was ahead of the local times at places of longitudes $x, x', x''$ etc. by the differences $a, a', a''$ etc. I assume that the readings $\theta, \theta', \theta''$ etc. have already been reduced to a single place. Thus if the daily rate of the chronometer is $=u$, one would have the $n-1$ equations if the chronometer were perfect.

\[
a-\theta u-x = a'-\theta' u-x' = a''-\theta'' u-x'' = a'''-\theta''' u-x''' = \text{etc.}
\]

\begingroup\rightskip=0pt plus 1em
For these equations to suffice to determine the unknown quantities $u, x, x', x''$ etc., one of the quantities $x, x', x''$ etc. is taken as given, and it is assumed that observations were made twice at at least one place, so that two of the quantities $x, x', x''$ etc. are identical. If no more than two are identical, the problem is fully determined. Otherwise it is overdetermined, and the unknown quantities must then be chosen so that the $n-1$ equations\par\endgroup

\[
\begin{aligned}
0 &= a-a'+(\theta'-\theta)u-x+x'\\
0 &= a'-a''+(\theta''-\theta')u-x'+x''\\
0 &= a''-a'''+(\theta'''-\theta'')u-x''+x''' \quad \text{etc.}
\end{aligned}
\]
\noindent are satisfied as closely as possible, since the imperfections that all chronometers invariably have will make it impossible to satisfy them all exactly. But these equations clearly must not be given equal weights, for the quantities

\[
\begin{aligned}
&a-a'+(\theta'-\theta)u-x+x'\\
&a'-a''+(\theta''-\theta')u-x'+x''
\end{aligned}
\]
\noindent express only the sums of all deviations from the mean rate of the chronometer during the intervals $\theta'-\theta$, $\theta''-\theta'$, and so on. For a good chronometer, to which one can indeed assign a mean rate free of gradually increasing changes in one direction, the mean value to be expected for such a sum must be taken as proportional to the square root of the interval.

\Needspace{3\baselineskip}
\SourceLeaf{461}
Thus, in treating the equations above by the method of least squares, one must assign them unequal weights inversely proportional to the intervals $\theta'-\theta$, $\theta''-\theta'$, $\theta'''-\theta''$, etc.

Solving the equations then presents no difficulty and yields both the most plausible values of $u, x, x', x''$ etc. and their relative reliability. I add a few remarks on this.

1) If the first and last observations were made at the same place, the most plausible value of $u$ is exactly the one that follows from comparing those two observations alone. The calculation then becomes extraordinarily simple, because an easily proved theorem permits direct substitution of this most plausible value of $u$ into the equations or, equivalently, reduction of all the observed chronometer times to those of an imaginary chronometer whose lead would be $=0$.

2) If the equations are simply given the weights $\dfrac{1}{\theta'-\theta}$, $\dfrac{1}{\theta''-\theta'}$, the unit underlying the weights found for the final longitude determinations is the accuracy one would expect with \emph{this} chronometer in determining a difference in longitude after an interval of one day, with its rate known (provided that the times $\theta, \theta', \theta''$ are expressed in days). To compare the results of \emph{different} chronometers of unequal quality, however, a further factor must be included that depends on the quality of each chronometer. To find it, set the values of the quantities

\[
\begin{aligned}
&a-a'+(\theta'-\theta)u-x+x'\\
&a'-a''+(\theta''-\theta')u-x'+x''\\
&a''-a'''+(\theta'''-\theta'')u-x''+x''' \quad \text{etc.}
\end{aligned}
\]
\noindent after substituting the most plausible values found for $u, x, x', x''$ etc., $=\lambda, \lambda', \lambda''$ etc. and

\[
\frac{\lambda\lambda}{\theta'-\theta} + \frac{\lambda'\lambda'}{\theta''-\theta'} + \frac{\lambda''\lambda''}{\theta'''-\theta''} + \text{etc.} = S
\]

Let $\nu$ also be the total number of quantities that were unknown, and let $m=\BareRoot{\dfrac{S'}{n-\nu-1}}$. The factor specific to each chronometer is then proportional to $\dfrac{1}{mm}$ or $\dfrac{n-\nu-1}{S}$. One may regard $m$ as the mean deviation from the mean rate to be expected after an interval of one day.

3) The rules above apply to a chronometer whose rate shows no appreciable progressive change. If the opposite is the case and the series of observations is not excessively long, one may be content to assume a change in the daily rate proportional to time. One more unknown quantity must then be introduced, and the equations take the following form:

\[
\begin{aligned}
0 &= a-a'+(\theta'-\theta)u+(\theta'\theta'-\theta\theta)v-x+x'\\
0 &= a'-a''+(\theta''-\theta')u+(\theta''\theta''-\theta'\theta')v-x'+x''\\
0 &= a''-a'''+(\theta'''-\theta'')u+(\theta'''\theta'''-\theta''\theta'')v-x''+x'''\\
&\text{etc.}
\end{aligned}
\]

4) It would hardly be advisable to go further and introduce yet another unknown quantity and terms of the form $(\theta'^3-\theta^3)w$. I would prefer to exclude altogether, along with others, chronometers that show large, distinct changes in their mean rate that are themselves irregular. Their results are much less accurate, and it is much harder to express their \SourceLeaf{462}accuracy numerically for comparison. I therefore do not dwell here on the much more complicated theory of such cases, since the foregoing will suffice in the present case.

4) As for solving the equations by the method of least squares, it may be worth recalling that in most cases it is best to write each unknown quantity as the sum of a known part (as close an approximation as possible) and an unknown part (thus a very small one). This advice has been given repeatedly elsewhere, and the advantage of this method is evident in itself. Still, it seemed worth recalling because I see it is often forgotten, which makes numerical calculations unnecessarily difficult and mistakes more likely.

Of the 36 chronometers\HistoricalNote{462} I have calculated the results for the following 5.

\begin{center}
\begingroup\fontsize{8}{10}\selectfont\setlength{\tabcolsep}{2pt}
\begin{tabular}{ll|r|r|r|r|r|r}
 & & & No. 1 & No. 4 & \textsc{Breguet} 3056 & \textsc{Kessels} 1252 & \textsc{Barraud} 904 \\
\hline
Greenwich & June 30 & $3^{\rm u}22^{\rm m}$ & $-8^{\rm m}17^{\rm s}14$ & $+1^{\rm m}2^{\rm s}37$ &  &  &  \\
 & July 25 & $2\;15$ & $10\;44.39$ & $1\;32.15$ & $+30^{\rm m}59^{\rm s}75$ & $+50^{\rm m}29^{\rm s}31$ & $+48^{\rm m}29^{\rm s}20$ \\
 & 28 & $3\;13$ & $11\;0.69$ & $1\;36.96$ & $30\;50.07$ & $50\;39.69$ & $48\;40.24$ \\
 & Aug. 2 & $1\;15$ & $11\;28.48$ & $1\;44.44$ & $30\;31.78$ & $50\;52.14$ & $48\;58.87$ \\
 & 17 & $10\;28$ & $12\;59.40$ & $2\;6.24$ & $29\;35.69$ & $51\;38.66$ & $49\;57.83$ \\
 & 25 & $7\;27$ & $13\;47.98$ & $2\;15.84$ & $29\;10.48$ & $52\;2.45$ & $50\;27.15$ \\
 & Sept. 10 & $7\;40$ & $15\;24.47$ & $2\;40.36$ &  &  &  \\
Helgoland & July 3 & $3\;40$ & $-40\;8.00$ & $-30\;26.84$ &  &  &  \\
 & 22 & $12\;40$ & $42\;2.02$ & $30\;3.89$ & $-0\;20.34$ & $+18\;48.39$ & $+16\;47.39$ \\
 & Aug. 5 & $1\;48$ & $43\;18.11$ & $29\;43.35$ & $1\;10.24$ & $19\;26.77$ & $17\;37.51$ \\
 & 11 & $13\;9$ & $43\;35.77$ & $29\;33.43$ & $1\;32.75$ & $19\;47.22$ & $18\;1.30$ \\
 & 30 & $19\;30$ & $45\;53.08$ & $29\;7.96$ & $2\;40.67$ & $20\;47.68$ & $19\;17.03$ \\
 & Sept. 6 & $3\;6$ & $46\;31.56$ & $28\;58.94$ & $3\;4.55$ & $21\;6.56$ & $19\;43.80$ \\
 & 7 & $8\;42$ & $46\;38.72$ & $28\;56.71$ &  &  &  \\
Altona & Aug. 6 & $5\;55$ & $-51\;38.95$ & $-37\;55.76$ & $-9\;28.50$ & $+11\;16.25$ & $+9\;28.48$ \\
 & 9 & $12\;35$ & $51\;57.35$ & $37\;50.03$ & $9\;38.81$ & $11\;27.76$ & $9\;40.30$ \\
 & 31 & $9\;57$ & $54\;10.33$ & $37\;21.30$ & $10\;56.68$ & $12\;35.96$ & $11\;5.92$ \\
 & Sept. 4 & $22\;12$ & $54\;39.16$ & $37\;15.21$ & $11\;15.36$ & $12\;48.10$ & $11\;24.49$ \\
Bremen & Aug. 13 & $0\;2$ & $-47\;50.65$ & $-33\;16.49$ & $-5\;23.37$ & $+16\;5.83$ & $+14\;21.86$ \\
\end{tabular}
\endgroup
\end{center}

I give the calculation for \emph{Breguet} 3056 as a check. Let the longitude of Helgoland be $=0$, that of Greenwich $=-x$, and that of Altona $=+y$. I leave Bremen out here because only one observation was made there, so it offers no check in any event.

I reckon the times from the first comparison of the English chronometers (Greenwich, 30 June, $3^{\rm u}22^{\rm m}$).

\Needspace{11\baselineskip}
By substituting an imaginary chronometer with rate $=0$, I find its offset as follows.

\begin{center}
\begin{tabular}{r|l||r|l}
$\theta$ & & $\theta$ & \\
\hline
22.4 & $+60^{\rm s}20$ & 42.4 & $+59^{\rm s}88$ \\
25.0 & $+1949.60-x$ & 48.3 & $+1949.60-x$ \\
28.0 & $+1950.87-x$ & 56.2 & $+1952.74-x$ \\
32.9 & $+1950.29-x$ & 61.6 & $+61.32$ \\
35.9 & $+59.08$ & 62.2 & $-432.53+y$ \\
37.1 & $-434.98+y$ & 66.8 & $-434.98+y$ \\
40.4 & $-433.49+y$ & 68.0 & $+60.19$ \\
\end{tabular}
\end{center}

\SourceLeaf{463}
In this case the equations above do not mix $x$ and $y$ at all, which makes the rest of the calculation even easier. We have four determinations of $x$.

\[
\begin{array}{ll}
+1889^{\rm s}40 & \text{Weight } \dfrac{1}{2.6}=0.38 \\[7pt]
+1891.21 & \dfrac{1}{3.0}=0.33 \\[7pt]
+1889.78 & \dfrac{1}{5.9}=0.17 \\[7pt]
+1891.42 & \dfrac{1}{3.4}=0.19 \\[7pt]
\hline
\text{Thus } x=+1890^{\rm s}36 \quad \text{Weight} & =1.07
\end{array}
\]

In the same way we find

\[
y=+494^{\rm s}12 \quad \text{Weight } =3.83
\]
\noindent Substituting these values gives the offset of the imaginary chronometer relative to Helgoland time

\begin{center}
\begingroup\small\renewcommand{\arraystretch}{1.25}
\begin{tabular}{r|r r@{\qquad}r|r r}
$\theta$ & & $\lambda$ & $\theta$ & & $\lambda$ \\
\cline{1-1}\cline{3-4}\cline{6-6}
$22.4$ & $+60^{\rm s}20$ & \raisebox{-7pt}{$-0^{\rm s}96$} & $42.4$ & $+59^{\rm s}88$ & \raisebox{-7pt}{$-0^{\rm s}62$} \\
$25.0$ & $59.24$ & \raisebox{-7pt}{$+1.27$} & $48.3$ & $59.24$ & \raisebox{-7pt}{$+3.14$} \\
$28.0$ & $60.51$ & \raisebox{-7pt}{$-0.58$} & $56.2$ & $62.38$ & \raisebox{-7pt}{$-1.06$} \\
$32.9$ & $59.93$ & \raisebox{-7pt}{$-0.85$} & $61.6$ & $61.32$ & \raisebox{-7pt}{$+0.27$} \\
$35.9$ & $59.08$ & \raisebox{-7pt}{$+0.06$} & $62.2$ & $61.59$ & \raisebox{-7pt}{$-2.45$} \\
$37.1$ & $59.14$ & \raisebox{-7pt}{$+1.47$} & $66.8$ & $59.14$ & \raisebox{-7pt}{$+1.05$} \\
$40.4$ & $60.63$ & \raisebox{-7pt}{$-0.75$} & $68.0$ & $60.19$ &  \\
$42.4$ & $59.88$ &  &  &  &  \\
\end{tabular}
\endgroup
\end{center}

Thus $S=6.00$, $m=\BareRoot{\dfrac{6.00}{13-3}}$. This gives the mean error to be expected for $x$: $0^{\rm s}75$, and for $y$: $0^{\rm s}40$

\Needspace{19\baselineskip}
The full set of 5 chronometers for which I calculated the results gives

\begin{center}
\begin{tabular}{ll|l|l}
 & & E.\,med. & Weight \\
\hline
Breguet & $x=1890^{\rm s}36$ & $0^{\rm s}75$ & 1.78 \\
Kessels & $1893.29$ & $0.67$ & 2.23 \\
Barraud & $1892.32$ & $0.49$ & 4.16 \\
Eng. 1 & $1892.39$ & $0.43$ & 5.41 \\
---\quad 4 & $1892.52$ & $0.35$ & 8.16 \\
\hline
Mean & $x=1892.35$ & & 21.74 \\
\end{tabular}
\end{center}

\begin{center}
\begin{tabular}{ll|l|l}
Breguet & $y=494.12$ & $0.40$ & 6.25 \\
Kessels & $493.89$ & $0.36$ & 7.72 \\
Barraud & $493.67$ & $0.26$ & 14.79 \\
Eng. 1 & $493.98$ & $0.29$ & 11.89 \\
---\quad 4 & $494.16$ & $0.24$ & 17.36 \\
\hline
 & $y=493.96$ & & 58.01 \\
\end{tabular}
\end{center}

In the last column, headed Weight, I have entered $\dfrac{1}{\text{Sq. (E.\,m.)}}$, taking as the unit the accuracy for which the mean error to be expected is $=1^{\rm s}$. For Altona, for example, the mean error to be expected is thus $=\dfrac{1^{\rm s}}{\BareRoot{58.01}}=0^{\rm s}13$. It would be more advisable, however, to regard the numbers in the last column merely as relative values and to infer the absolute accuracy from the differences between the values of $x$ and $y$ obtained from the individual chronometers and the final results. Even then the accuracy of the final result \SourceLeaf{464}will still appear somewhat greater than it really is, since the time determinations at Greenwich, Helgoland and Altona are not \emph{absolutely} accurate. Clearly, however many chronometers there might be, errors arising from that source must still affect the final results.

The longitude of Bremen can be determined as follows. If we take the longitude to be $=z$ east of Helgoland, the comparison of the \textsc{Breguet} chronometer gives

\[
-165^{\rm s}52+z
\]


\begin{center}\small
\begin{tabular}{p{0.42\linewidth} r l r}
Thus from the preceding comparison & $z=225.40$ & Weight & $\dfrac{1}{1.4}=0.7$ \\[5pt]
from the following comparison & $z=224.76$ & & $\dfrac{1}{4.5}=0.2$ \\[8pt]
\cline{2-2}\cline{4-4}
 & $225.24$ & & $0.9$ \\
\end{tabular}
\end{center}

\Needspace{12\baselineskip}
The weight $0.9$ must still be multiplied by $\dfrac{10}{6.00}$. The 5 chronometers thus give

\begin{center}
\begin{tabular}{l|l|l}
 & & Weight \\
\hline
Breguet & $225.24$ & $1.5$ \\
Kessels & $225.84$ & $1.9$ \\
Barraud & $225.39$ & $3.6$ \\
Eng. 1 & $226.04$ & $2.9$ \\
---\quad 4 & $224.86$ & $4.3$ \\
\hline
 & $225.42$ & $14.2$ \\
\end{tabular}
\end{center}

But the longitude of Bremen, which according to this calculation comes out $268^{\rm s}54$ west of Altona, naturally still depends on the time determination in Bremen, and this difference seems several seconds too small. According to my triangles, the Ansgariusthurm is $273^{\rm s}51$ in time west of Göttingen, so \textsc{Olbers}’s observatory is $271^{\rm s}9$.
